\documentclass[11pt,UTF8,fontset=fandol]{ctexart}
\usepackage{amsmath,amssymb,amsfonts}
\usepackage{graphicx,float,caption,subcaption}
\usepackage{booktabs,array,multirow,makecell}
\usepackage{geometry,enumitem,xcolor}
\usepackage[hidelinks]{hyperref}
\usepackage{url}
\geometry{margin=2.3cm}
\graphicspath{{/Users/yye/project/论文/assets/searchr1/}}
\captionsetup{font=small,labelfont=bf}
\setlength{\parskip}{0.4em}

% 自定义：模型输出标签盒子，便于展示 <think>/<search> 等结构
\newcommand{\tg}[1]{\texttt{<#1>}}

\title{\textbf{Search-R1：用强化学习训练大模型“边推理边自主调用搜索引擎”}}
\author{Bowen Jin, Hansi Zeng, Zhenrui Yue, Jinsung Yoon, Sercan Ö. Arık,\\
Dong Wang, Hamed Zamani, Jiawei Han\\[0.3em]
\small University of Illinois at Urbana-Champaign（伊利诺伊大学厄巴纳-香槟分校）\\
\small University of Massachusetts Amherst（马萨诸塞大学阿默斯特分校）\\
\small Google Cloud AI Research（谷歌云 AI 研究院）}
\date{}

\begin{document}
\maketitle

\begin{center}
\fbox{\parbox{0.95\linewidth}{\small
\textbf{论文信息}\quad 会议/年份：COLM 2025（arXiv:2503.09516）\quad|\quad 引用数：1068（Semantic Scholar，2026-06-07，本领域最高）\\
\textbf{代码}：\url{https://github.com/PeterGriffinJin/Search-R1}\\[0.3em]
\textbf{一句话定位}：这篇论文要解决的问题是——\emph{怎么让大模型在自己思考的过程中，像人一样“想到不懂就去搜一下、读完再接着想、必要时再搜”，而且这种“边想边搜”的本事是被强化学习\textbf{训练出来}的，不是靠人写提示词硬凑出来的}。它是“搜索类/深度研究类（Deep-Research）智能体用强化学习训练”这条技术路线的奠基之作。
}}
\end{center}

\begin{abstract}
\noindent\textbf{摘要（全文翻译）}\quad 对大语言模型（large language models，LLMs）来说，高效地获取外部知识与最新信息，对于有效的推理和文本生成至关重要。直接用提示词（prompting）让具备推理能力的先进 LLM 在推理（inference）时去使用搜索引擎，往往效果并不理想，因为模型可能并不真正掌握“如何与搜索引擎进行最优交互”的能力。本文提出 Search-R1，它是对“面向推理的强化学习（reinforcement learning，RL）框架”的一种扩展：让 LLM 学会在逐步推理的过程中，配合实时检索，\emph{自主地生成（多次）搜索查询}。Search-R1 通过带多轮搜索交互的方式来优化 LLM 的推理轨迹（trajectory），并借助\emph{对检索回来的 token 做损失掩码（retrieved token masking）}来实现稳定的 RL 训练，同时采用一个简单的、基于最终结果的奖励函数（outcome-based reward）。在 7 个问答（question-answering）数据集上的实验表明，在相同实验设置下，Search-R1 相对各类 RAG 基线分别提升了 24\%（Qwen2.5-7B）和 20\%（Qwen2.5-3B）。本文还就 RL 优化方法、LLM 模型选择，以及检索增强推理中“回复长度的动态变化”给出了实证洞察。代码与模型检查点（checkpoints）已开源：\url{https://github.com/PeterGriffinJin/Search-R1}。
\end{abstract}

%==================================================================
\section{小白知识卡：读这篇前要懂的几个词}
\label{sec:cards}

下面这些概念在本文反复出现，先用大白话过一遍，后面读正文就轻松了。

\begin{itemize}[leftmargin=1.4em]
\item \textbf{大语言模型（LLM）}：就是会“接话”的模型（如 GPT、Qwen）。它根据已经看到的文字，一个 token 一个 token 地往下预测、生成内容。

\item \textbf{RAG（检索增强生成，retrieval-augmented generation）}：传统做法——先拿用户问题去知识库里“搜一把”，把搜到的几段文字塞进模型的输入里，模型再据此作答。缺点是\emph{只搜一次}、而且模型本身并没有被训练“怎么搜更聪明”。

\item \textbf{把搜索当工具（search-as-a-tool）}：让模型在思考过程中\emph{主动}发起搜索（像 ReAct、IRCoT 那样）。但若只靠写提示词，模型常常用不好；若靠监督微调（SFT），又得人工标注大量“正确的搜索轨迹”，成本极高。

\item \textbf{强化学习（RL）}：一种“试错+给分”的训练范式。模型（这里叫\emph{策略 policy}）去做一连串决策，环境给它一个奖励（reward）分数，模型据此调整自己、争取拿更高分。它最大的好处是：\emph{只告诉模型“最终对没对”，模型自己就能摸索出复杂的中间策略}。

\item \textbf{PPO（近端策略优化，proximal policy optimization）}：最常用的一种 RL 算法。它需要额外训练一个“裁判”（叫\emph{价值模型/critic}）来估计“当前局面值多少分”，训练更稳但更重。

\item \textbf{GRPO（组相对策略优化，group relative policy optimization）}：DeepSeek 提出的简化版 RL。它不另外训裁判，而是对同一个问题\emph{采样一组答案}，用这组答案的平均分当基准，谁高于平均就鼓励、谁低于平均就抑制。更轻、收敛更快，但有时会“崩”。

\item \textbf{RLVR（带可验证奖励的强化学习，RL with verifiable rewards）}：奖励不靠另一个神经网络打分，而是\emph{用规则直接判对错}（比如答案字符串是否精确匹配标准答案）。本文用的就是这种，简单又不容易被模型“钻空子”。

\item \textbf{DeepSeek-R1 / R1-Zero}：用纯 RL（只给结果奖励）就让模型自己学会长链推理、自我验证、自我纠错的代表作。\textbf{本文可以理解成它的“加搜索版”}：R1 让模型“在脑子里想”，Search-R1 让模型“想不通就上网查”。

\item \textbf{token 掩码（token masking）}：训练时把某些 token 从“算账”里排除。本文的关键技巧——\emph{搜索引擎返回的那段文字不是模型自己写的，所以训练时不让模型为这段内容的概率“负责/背锅”}，否则训练会乱套。
\end{itemize}

%==================================================================
\section{引言在讲什么（对应原文第 1 节）}

\textbf{背景痛点。}\quad LLM 很能“说”，但有两个老毛病：（1）做复杂推理容易出错；（2）缺乏\emph{最新的、领域专门的}外部知识，还爱“一本正经地胡说”（即幻觉，hallucination）。要治这两个病，就得让模型既会推理、又会高效地用搜索引擎调外部信息。

\textbf{现有两条路都不够好。}\quad 论文把已有做法分成两类：
\begin{itemize}[leftmargin=1.4em]
\item \emph{RAG}：基于输入去检索段落、拼进上下文再生成。哪怕有人改成“多轮多次检索”，模型也没在训练中\emph{学过}怎么和搜索引擎打交道，所以是次优的。
\item \emph{把搜索当工具}：用提示词驱动（泛化差，没见过的任务就抓瞎）或用监督微调（要海量高质量标注轨迹，难以规模化）。更要命的是——\emph{搜索这个动作本身是“不可微”的}（你没法对“去网上搜一下”这个离散操作求梯度），所以端到端的梯度下降训不动它。
\end{itemize}

\textbf{RL 是突破口，但有三道坎。}\quad 既然 OpenAI-o1、DeepSeek-R1 证明了“只给结果奖励、用 RL”就能逼出强推理能力，那能不能把搜索也用 RL 训出来？作者指出三个核心挑战：
\begin{enumerate}[leftmargin=1.6em]
\item \textbf{RL 框架与稳定性}：怎么把搜索引擎\emph{塞进}RL 流程，尤其是检索回来的上下文混进序列后，怎么保证训练稳定？
\item \textbf{多轮交错的“推理-搜索”}：理想情况下，模型应能反复地“想—搜—再想—再搜”，并根据问题难度动态调整检索策略。
\item \textbf{奖励设计}：到底用不用复杂的“过程奖励”？还是说一个简单的“只看最终答案对错”的奖励就够了？
\end{enumerate}

\textbf{本文的答案——Search-R1 的三个关键创新：}
\begin{enumerate}[leftmargin=1.6em]
\item \emph{把搜索引擎建模为环境的一部分}，于是采样出的轨迹就是“模型生成的 token”和“搜索返回的 token”交错在一起的序列。框架同时兼容 PPO 与 GRPO，并用\textbf{检索 token 掩码}保证训练稳定。
\item \emph{支持多轮检索与推理}：用 \tg{search}…\tg{/search} 触发搜索，检索结果包在 \tg{information}…\tg{/information} 里，推理过程包在 \tg{think}…\tg{/think} 里，最终答案放在 \tg{answer}…\tg{/answer} 里，形成结构化的迭代决策。
\item \emph{只用极简的“结果奖励”}（答对就给分），不搞复杂的过程奖励。实验证明这个最小化设计就够用。
\end{enumerate}

可以把 Search-R1 看成 DeepSeek-R1-Zero 的扩展：后者主攻“用模型自身参数里的知识做推理”，前者引入“搜索增强的 RL 训练”，让决策能由外部检索来驱动。

%==================================================================
\section{相关工作速览（对应原文第 2 节）}

\textbf{LLM 与检索。}\quad 给 LLM 接外部信息主要两条路：RAG（检索一次→顺序生成，问题是可能检到无关内容、上下文不够用）和把搜索当工具（IRCoT、ReAct 用提示词驱动迭代推理+搜索；Toolformer 用监督微调增强用工具的能力，但依赖难以规模化的高质量标注轨迹）。最近的工作（DeepSeek-R1）显示：\emph{只靠结果奖励，RL 就能逼出高级推理}，但它在“调用搜索引擎”这个场景下还没被充分探索——这正是本文的切入点。

\textbf{LLM 与强化学习。}\quad RL 进入 LLM 调优是从 RLHF（基于人类反馈的强化学习）开始的：先用人类偏好数据训一个奖励模型，再用它（通常是 PPO）去调策略 LLM。但 PPO 工程上较重。于是出现了 DPO（直接偏好优化）、SimPO 等\emph{直接优化}方法，省事但有“离策略（off-policy）”问题，性能不总能赶上纯 RL。另一支是 GRPO（去掉 critic，用一组样本的平均分当基准）和 RLOO（简化版 REINFORCE）。但把 RL 用到“LLM 驱动的搜索交互+推理”上，此前几乎是空白。

%==================================================================
\section{方法：Search-R1 怎么做（对应原文第 3 节）}

这一节是全文核心，分四块讲：（1）把搜索引擎接进 RL 的目标函数；（2）多轮搜索的生成流程；（3）训练用的模板；（4）奖励设计。

\subsection{3.1 把搜索引擎接进强化学习}

\textbf{优化目标（公式 1）。}\quad 作者把“带搜索引擎 $R$”的 RL 目标写成：
\begin{equation*}
\max_{\pi_\theta}\ \mathbb{E}_{x\sim D,\, y\sim \pi_\theta(\cdot\,|\,x;R)}\big[\, r_\phi(x,y)\,\big]\ -\ \beta\, D_{\mathrm{KL}}\!\big[\pi_\theta(y\,|\,x;R)\ \|\ \pi_{\mathrm{ref}}(y\,|\,x;R)\big].
\tag{1}
\end{equation*}
\textbf{逐符号拆解}：
\begin{itemize}[leftmargin=1.4em]
\item $\pi_\theta$：要训练的\emph{策略模型}（就是我们的 LLM）；$\pi_{\mathrm{ref}}$：\emph{参考模型}（训练前的初始模型，当“锚点”用，防止训歪）。
\item $x\sim D$：从数据集 $D$ 里抽一个输入（问题）；$r_\phi(x,y)$：\emph{奖励函数}，给“问题 $x$ + 回答 $y$”打分。
\item $y\sim \pi_\theta(\cdot\,|\,x;R)$：注意这里的 $y$ 不再是模型一口气写完的，而是\textbf{“推理 token”和“搜索引擎检索结果”交错}组成的序列——这就是和普通 RL 最大的区别。作者把它记作 $\pi_\theta(\cdot\,|\,x)\otimes R$，$\otimes$ 表示“检索与推理交错进行”。
\item $D_{\mathrm{KL}}[\cdot\|\cdot]$：KL 散度，衡量新策略和参考模型差多远；$\beta$ 是它的权重。这一项是“别跑太远”的缰绳，保证训练稳定。
\end{itemize}
直觉：在“答对就给分”的驱动下，模型自己去摸索“什么时候该搜、搜什么、怎么用搜到的东西”，同时别离初始模型太远以免崩。框架建立在两大经典策略梯度方法 PPO 和 GRPO 之上。

\textbf{关键技巧：对检索 token 做损失掩码（Loss Masking for Retrieved Tokens）。}\quad 这是本文最重要的工程洞见，务必理解：
\begin{quote}\small
在 PPO/GRPO 里，token 级的损失是\emph{对整条轨迹}逐 token 计算的。而 Search-R1 的轨迹里，既有“模型自己生成的 token”，也有“从外部段落检索回来的 token”。\textbf{优化模型自己写的 token}（怎么发起搜索、怎么推理）是我们想要的；但如果\textbf{把同样的优化也套到检索回来的 token 上}，就会导致“意料之外的学习动态”——说白了，搜索引擎返回什么内容\emph{不是模型能控制的}，硬要模型去“拟合/对这段内容负责”就会把训练带乱。
\end{quote}
解决办法：引入一个掩码 $I(y_t)$——若 $y_t$ 是模型生成的 token 则 $I(y_t)=1$，若是检索回来的 token 则 $I(y_t)=0$；策略梯度\emph{只在模型生成的 token 上计算}。这既稳定了训练，又保留了“搜索增强生成”的灵活性。\textbf{这是 Search-R1 能稳定训练的命门。}

\textbf{PPO 版（公式 2）。}\quad PPO 是常用的 actor-critic 方法，其（带搜索的）目标为：
\begin{equation*}
J_{\mathrm{PPO}}(\theta)=\mathbb{E}_{x\sim D,\,y\sim\pi_{\mathrm{old}}}\!\left[\frac{1}{\sum_{t=1}^{|y|}I(y_t)}\!\!\sum_{t=1:\,I(y_t)=1}^{|y|}\!\!\min\!\Big(\tfrac{\pi_\theta(y_t|x,y_{<t};R)}{\pi_{\mathrm{old}}(y_t|x,y_{<t};R)}A_t,\ \mathrm{clip}\big(\tfrac{\pi_\theta(y_t|x,y_{<t};R)}{\pi_{\mathrm{old}}(y_t|x,y_{<t};R)},1{-}\epsilon,1{+}\epsilon\big)A_t\Big)\right].
\tag{2}
\end{equation*}
\textbf{怎么读}：$\pi_\theta$ 是当前策略、$\pi_{\mathrm{old}}$ 是上一步的旧策略，二者之比是\emph{重要性采样比}（衡量“同一个 token 在新旧策略下概率变化多少”）。$A_t$ 是\emph{优势（advantage）}，表示“这一步比平均水平好多少”，由 GAE（广义优势估计）基于未来奖励和一个学到的价值函数 $V_\phi$ 算出。$\mathrm{clip}(\cdot,1{-}\epsilon,1{+}\epsilon)$ 把这个比值\emph{裁剪}在 $[1{-}\epsilon,1{+}\epsilon]$，防止单步更新过猛——这正是 PPO 稳定的来源。注意求和写的是 $t=1:I(y_t)=1$、分母是 $\sum I(y_t)$：\textbf{只对“模型自己写的 token”求和、求平均}，检索 token 被掩掉了。

\textbf{GRPO 版（公式 3）。}\quad 为了不用额外的价值函数、提升稳定性，GRPO 对每个问题 $x$ 采样一组回答 $\{y_1,\dots,y_G\}$，用\emph{组内相对奖励}来算优势 $\hat A_{i,t}$（谁比这一组的平均好就鼓励）：
\begin{equation*}
\begin{aligned}
J_{\mathrm{GRPO}}(\theta)=\mathbb{E}\Big[\tfrac{1}{G}\sum_{i=1}^{G}\tfrac{1}{\sum_t I(y_{i,t})}\!\!\sum_{t:\,I(y_{i,t})=1}\!\!\Big(&\min\big(\tfrac{\pi_\theta(y_{i,t}|\cdot)}{\pi_{\mathrm{old}}(y_{i,t}|\cdot)}\hat A_{i,t},\ \mathrm{clip}(\tfrac{\pi_\theta(y_{i,t}|\cdot)}{\pi_{\mathrm{old}}(y_{i,t}|\cdot)},1{-}\epsilon,1{+}\epsilon)\hat A_{i,t}\big)\Big)\\[-0.2em]
&-\ \beta\, D_{\mathrm{KL}}[\pi_\theta\|\pi_{\mathrm{ref}}]\Big].
\end{aligned}
\tag{3}
\end{equation*}
\textbf{和 PPO 的两点不同}：（1）优势 $\hat A_{i,t}$ 用\emph{组内相对奖励}算，不需要 critic；（2）KL 惩罚不是塞进奖励里，而是\emph{直接加在损失上}。同样地，算 KL 损失时\textbf{也对检索 token 做掩码}。

\subsection{3.2 多轮搜索引擎调用的生成流程}

模型的回答不是一次成型，而是\emph{在“文本生成”和“查搜索引擎”之间反复切换}。系统提示告诉模型：需要检索时，把查询写在 \tg{search}…\tg{/search} 之间；系统一旦在生成里检测到这对标签，就提取查询、去搜、把结果包在 \tg{information}…\tg{/information} 里追加到序列后面，作为下一步生成的额外上下文。如此往复，直到：（1）达到最大动作次数，或（2）模型给出最终答案（包在 \tg{answer}…\tg{/answer} 里）。完整流程见原文算法 1（Algorithm 1），核心循环用大白话说就是：

\begin{quote}\small
\textbf{算法 1（多轮搜索的回复生成，大白话版）}：给定问题 $x$、策略模型、搜索引擎、最大动作预算 $B$。\\
重复至多 $B$ 次：①让模型往下生成 token，直到吐出 \tg{/search}、\tg{/answer} 或结束符就停；②如果这一段里出现了 \tg{search}…\tg{/search}，就把查询抠出来、调搜索引擎拿结果 $d$，并以 \tg{information}$d$\tg{/information} 的形式插回序列；③如果出现了 \tg{answer}…\tg{/answer}，就返回最终答案；④否则（动作不合法）追加一句“My action is not correct. Let me rethink.”逼模型重想。
\end{quote}

\subsection{3.3 训练模板（对应原文表 1）}

为了引导初始模型按格式输出，作者用了一个极简模板，把输出结构化成三部分循环：先\emph{推理}，再\emph{发起搜索}，最后\emph{给答案}。\textbf{关键是：模板只规定“格式”，不规定“内容”}——不强迫模型必须反思、必须搜索、也不偏向某种解题套路，这样才能\emph{原汁原味地观察}RL 过程中模型自发的学习动态。模板原文（保留英文，便于复现）大意是：

\begin{quote}\small\itshape
回答给定的问题。每次拿到新信息，你必须先在 \tg{think} 和 \tg{/think} 之间进行推理。推理之后，若发现自己缺少某些知识，可以用 \tg{search} 查询 \tg{/search} 来调用搜索引擎，它会把检索到的 top 结果返回在 \tg{information} 和 \tg{/information} 之间。你想搜多少次都行。若发现不再需要外部知识，就可以直接把答案写在 \tg{answer} 和 \tg{/answer} 之间，不必详细解释。例如 \tg{answer} xxx \tg{/answer}。问题：\{question\}。
\end{quote}

\subsection{3.4 奖励设计（公式 4）}

奖励是 RL 的主要训练信号。Search-R1 采用\emph{纯基于规则、只看最终结果}的奖励系统，用来判断模型答得对不对。例如事实类推理任务，可以用精确字符串匹配（exact match，EM）：
\begin{equation*}
r_\phi(x,y)=\mathrm{EM}(a_{\mathrm{pred}},\, a_{\mathrm{gold}}),
\tag{4}
\end{equation*}
其中 $a_{\mathrm{pred}}$ 是从回答 $y$ 里抽出的最终答案，$a_{\mathrm{gold}}$ 是标准答案。两点设计取舍值得注意：（1）\textbf{不加格式奖励}——因为训出来的模型本身就很守格式（这点和 DeepSeek-R1 不同，R1 用了格式奖励）；（2）\textbf{不训神经网络奖励模型}——因为大规模 RL 下 LLM 对奖励形式很敏感、易被“钻空子”，而且再训一个奖励模型又费算力又复杂。\emph{一句话：奖励越简单、越“可验证”，越稳。}

%==================================================================
\section{方法图解：一张图看懂训练流程}

图~\ref{fig:arch} 是全文的“总览图”，展示了 Search-R1 分别用 PPO（上半部分）和 GRPO（下半部分）训练时，数据是怎么流动的。

\begin{figure}[H]\centering
  \includegraphics[width=0.97\linewidth]{cand_p04_k00.pdf}
  \caption{用搜索引擎进行 PPO 与 GRPO 训练的示意图（Search-R1）。在 rollout（轨迹采样）阶段，LLM 可以与搜索引擎进行多轮交互。图例：黄色=被训练的模型，绿色=冻结不动的模型，蓝色=搜索引擎。}
  \label{fig:arch}
\end{figure}

\textbf{图里在说什么、数据怎么流动？}

\begin{itemize}[leftmargin=1.4em]
\item \textbf{Rollout Module（轨迹采样模块，左侧灰框）}：这是和别的 RL 框架最不一样的地方。问题 $q$ 进来后，\emph{Policy LLM（策略模型，黄色，要训练）}和 \emph{Search Engine（搜索引擎，蓝色）}\textbf{反复来回交互}——模型生成查询→搜索引擎返回结果→模型接着生成……最终产出一条“推理与检索交错”的完整轨迹 $o$（GRPO 下是一组 $o_1,\dots,o_G$）。

\item \textbf{上半部分（PPO）}：轨迹 $o$ 分别送给三个模型——\emph{Value LLM（价值模型，黄色，要训练）}输出价值 $v$；\emph{Reward Model（奖励模型）}给出奖励 $r$；\emph{Reference LLM（参考模型，绿色，冻结）}用来算 KL（图中 $\oplus$ 表示把 KL 惩罚并入）。然后 $v$ 和 $r$ 一起喂给 \textbf{GAE} 算出优势 $A$，再用 $A$ 去更新策略。注意这里需要一个额外的 Value LLM，所以 PPO“更重”。

\item \textbf{下半部分（GRPO）}：对同一个 $q$ 采样一\emph{组}轨迹 $o_1,\dots,o_G$，每条各自拿到奖励 $r_1,\dots,r_G$，送进 \textbf{Group Computation（组内计算）}，用“组内相对好坏”直接算出各自的优势 $A_1,\dots,A_G$。\textbf{没有 Value LLM、没有 GAE}，所以 GRPO“更轻、收敛更快”，但也更容易不稳定（后面实验会看到它训久了可能“崩”）。Reference LLM 仍用于算 KL。
\end{itemize}

一句话总结这张图：\emph{两条路线的唯一公共关键就是左边那个“模型↔搜索引擎来回交互”的 Rollout Module；右边只是 PPO 和 GRPO 各自不同的“怎么把奖励折算成更新信号”的算法差异}。

%==================================================================
\section{实验解读（对应原文第 4 节）}

\subsection{用了哪些数据集、对比哪些基线}

\textbf{7 个问答数据集}（涵盖单跳与多跳，既测“训练域内”也测“域外泛化”）：
\begin{itemize}[leftmargin=1.4em]
\item \emph{通用问答（General QA）}：NQ（Natural Questions）、TriviaQA、PopQA。
\item \emph{多跳问答（Multi-Hop QA，需要分步推理、串联多条线索）}：HotpotQA、2WikiMultiHopQA、Musique、Bamboogle。
\end{itemize}
其中 \textbf{NQ 和 HotpotQA 是“训练域内”}（in-domain，记 $\dagger$，训练数据就是这俩合并而成），其余 5 个是\textbf{“训练域外”}（out-of-domain，记 $\star$，从没在训练里见过，用来检验泛化）。

\textbf{对比的基线（baselines）三大类}：（1）\emph{不检索}：直接推理（Direct Inference）、思维链（Chain-of-Thought, CoT）；（2）\emph{带检索}：RAG、IRCoT、Search-o1；（3）\emph{微调类}：SFT（监督微调）、R1（用同样数据做 RL 但\emph{不带搜索}）、带搜索的拒绝采样（rejection sampling，挑出答对的轨迹再去微调）。

\textbf{公平性设置}：所有方法用\emph{同一个}检索器（E5）、\emph{同一个}知识库（2018 年维基百科）、\emph{同样}检索 3 段（top-3）、\emph{同样}的训练数据和\emph{同样}的基座模型（Qwen2.5-3B/7B 的 Base/Instruct 版）。评测指标用 EM（精确匹配）。默认 RL 算法是 PPO。

\subsection{主结果：Search-R1 全面领先}

表~\ref{tab:main} 是核心结果（节选 7B 与 3B 的关键行，完整数据见原文表 2）。

\begin{table}[H]\centering
\caption{主结果（节选）。数值为 EM（越高越好），$\dagger$=训练域内、$\star$=训练域外；\textbf{加粗}为该规模下最佳。Avg.=七数据集平均。}
\label{tab:main}
\small
\begin{tabular}{lccccccc>{\bfseries}c}
\toprule
方法 & NQ$^\dagger$ & TriviaQA$^\star$ & PopQA$^\star$ & HotpotQA$^\dagger$ & 2wiki$^\star$ & Musique$^\star$ & Bamboogle$^\star$ & Avg. \\
\midrule
\multicolumn{9}{l}{\emph{Qwen2.5-7B}}\\
Direct Inference   & 0.134 & 0.408 & 0.140 & 0.183 & 0.250 & 0.031 & 0.120 & 0.181 \\
RAG                & 0.349 & 0.585 & 0.392 & 0.299 & 0.235 & 0.058 & 0.208 & 0.304 \\
R1-base（无搜索RL）  & 0.297 & 0.539 & 0.202 & 0.242 & 0.273 & 0.083 & 0.296 & 0.276 \\
Rejection Sampling & 0.360 & 0.592 & 0.380 & 0.331 & 0.296 & 0.123 & 0.355 & 0.348 \\
\textbf{Search-R1-base}     & \textbf{0.480} & \textbf{0.638} & \textbf{0.457} & \textbf{0.433} & 0.382 & \textbf{0.196} & \textbf{0.432} & \textbf{0.431} \\
\textbf{Search-R1-instruct} & 0.393 & 0.610 & 0.397 & 0.370 & \textbf{0.414} & 0.146 & 0.368 & 0.385 \\
\midrule
\multicolumn{9}{l}{\emph{Qwen2.5-3B}}\\
RAG                & 0.348 & 0.544 & 0.387 & 0.255 & 0.226 & 0.047 & 0.080 & 0.270 \\
R1-base（无搜索RL）  & 0.226 & 0.455 & 0.173 & 0.201 & 0.268 & 0.055 & 0.224 & 0.229 \\
\textbf{Search-R1-base}     & \textbf{0.406} & \textbf{0.587} & \textbf{0.435} & 0.284 & 0.273 & 0.049 & 0.088 & 0.303 \\
\textbf{Search-R1-instruct} & 0.341 & 0.545 & 0.378 & \textbf{0.324} & \textbf{0.319} & \textbf{0.103} & \textbf{0.264} & \textbf{0.325} \\
\bottomrule
\end{tabular}
\end{table}

\textbf{这些数字说明了什么（原文四点结论）：}
\begin{enumerate}[leftmargin=1.6em]
\item \textbf{Search-R1 稳定碾压强基线}：相对各类 RAG 基线，7B 平均提升 \emph{24\%}、3B 提升 \emph{20\%}（论文摘要口径；按“两个模型平均相对提升”另一口径为 41\% 和 20\%）。无论域内（NQ/HotpotQA）还是域外（其余 5 个）都成立——说明它\emph{学到的是通用的“会搜会想”能力，而非死记某几个数据集}。
\item \textbf{比“无搜索的 RL（R1）”更强}：把搜索引入推理，等于给模型开了“查资料”的外挂，自然更准。这也反过来说明：\emph{光靠模型脑子里的知识不够，接外部检索是真有用}。
\item \textbf{Base 和 Instruct 模型都吃这一套}：证明 DeepSeek-R1-Zero 那种“只用结果奖励的 RL”不仅能练纯推理，\emph{也能成功迁移到“推理+搜索”场景}。
\item \textbf{模型越大越会搜}：7B 相对第二名（RAG）的“性能差距”明显大于 3B——大模型更能学会“如何把搜索用好”。
\end{enumerate}

\subsection{消融与分析（对应原文第 5 节，几条最有用的）}

\textbf{（1）PPO vs. GRPO}（原文图 2a、表 3）：
\begin{itemize}[leftmargin=1.4em]
\item GRPO \emph{收敛更快}（因为 PPO 要先花几步“热身”那个 critic）；
\item 但 PPO \emph{更稳}——GRPO 训久了会出现“奖励崩塌（reward collapse）”，PPO 始终平稳；
\item 两者\emph{最终性能相当}。所以本文默认用 PPO（求稳）。
\end{itemize}

\textbf{（2）检索 token 掩码到底有多重要}（原文图 3、表 4）：这是验证那个“命门技巧”的关键消融。结果非常干脆——\emph{加掩码全面优于不加}：

\begin{table}[H]\centering
\caption{检索 token 损失掩码的消融（Qwen2.5-7B-base，PPO）。加掩码在全部 7 个数据集上一致更好。}
\label{tab:mask}
\small
\begin{tabular}{lccccccc>{\bfseries}c}
\toprule
方法 & NQ & TriviaQA & PopQA & HotpotQA & 2wiki & Musique & Bamboogle & Avg. \\
\midrule
Search-R1 \textbf{加}掩码   & 0.480 & 0.638 & 0.457 & 0.433 & 0.382 & 0.196 & 0.432 & 0.431 \\
Search-R1 不加掩码 & 0.388 & 0.567 & 0.391 & 0.325 & 0.321 & 0.108 & 0.304 & 0.343 \\
\bottomrule
\end{tabular}
\end{table}
平均分从 0.343 跳到 0.431——\textbf{“不让模型为搜索引擎返回的内容背锅”这一招，直接带来约 0.09 的平均提升}，是 Search-R1 能稳定见效的核心原因之一。

\textbf{（3）回复长度与搜索次数的动态}（原文图 2c、2d）：训练过程呈现有趣的“先降后升再稳”规律——
\begin{itemize}[leftmargin=1.4em]
\item \emph{前 100 步}：回复长度\emph{骤降}，奖励微升。此阶段模型在“删废话”，去掉啰嗦的填充词、适应任务。
\item \emph{100 步以后}：回复长度和奖励\emph{一起显著上升}。此时模型\emph{学会频繁调用搜索引擎}，检索段落让回复变长，而善用搜索结果让奖励大涨。
\item 与之对应，\emph{有效搜索次数随训练持续增多}——\textbf{模型是“自发”学会了多搜、反复搜的}，没人教它。
\end{itemize}

\textbf{（4）Base vs. Instruct}（原文图 2b）：Instruct 模型起点高、收敛快，但\emph{最终}两者性能很接近——说明 RL 有能力把差距补平，让 Base 模型也达到相当水平。

\textbf{（5）其它（附录）}：在 14B 模型上 Search-R1 依然全面领先且随规模继续涨（附录 C）；检索段数 top-$k$=3 综合最优（太少召回不够、太多引入噪声反而拖累，附录 G）；GRPO 的组大小（group size）越大收敛越快但越易崩，小组更稳、泛化更好（附录 H）。

\subsection{案例研究：模型“边想边搜、还会自我验证”}

原文附录用大量案例展示了训练后模型的行为。最经典的一例（对比 R1 与 Search-R1，原文表 9）：

\begin{quote}\small
\textbf{问题}：Curious 这款女士香水的代言歌手，出生在哪个城市和州？（标准答案：McComb, Mississippi）\\[0.3em]
\textbf{R1（无搜索，只靠脑补）}：\tg{think} 这位歌手是 Beyoncé，出生在休斯顿……\tg{/think} \tg{answer} Houston \tg{/answer} —— \textcolor{red}{答错了（连歌手都认错）}。\\[0.3em]
\textbf{Search-R1（边搜边想）}：先搜“Curious fragrance information”→得知是 Britney Spears 代言；再搜“Britney Spears birthplace”→得知出生地 McComb, Mississippi；\textbf{此时其实已经够了}，但模型\emph{又多搜了一次}“McComb, Mississippi location”来\textbf{自我验证}，最后 \tg{answer} McComb, Mississippi \tg{/answer} —— \textcolor{teal}{答对，且过程严谨}。
\end{quote}

作者从案例里总结出两点：（1）\textbf{交错的推理与检索能加深问题分析}——模型可以一边搜一边迭代地修正自己的推理；（2）\textbf{通过反复检索做自我验证（self-verification）}——哪怕信息已足够，模型仍会再搜一次确认，这和“无搜索的推理 RL（DeepSeek-R1）”里观察到的“验证驱动的推理”一脉相承。当然，附录也诚实地给了\emph{失败案例}：复杂问题分解不当、或被无关检索结果误导时，模型仍会答错（如把多跳问题想偏、被噪声段落带跑）。

%==================================================================
\section{核心 takeaways 与对 AI 应用开发的启示}

\textbf{这篇论文最值得记住的几点：}
\begin{enumerate}[leftmargin=1.6em]
\item \textbf{“会搜”可以被训练，而不是靠提示词硬凑}。把搜索引擎当成\emph{环境的一部分}，用 RL 端到端地训，模型就能自主学会“何时搜、搜什么、怎么用”。这是从“RAG/工具调用靠人设计”到“搜索行为靠模型自学”的范式转变。
\item \textbf{最关键的工程技巧是“检索 token 掩码”}：训练时\emph{不让模型为搜索引擎返回的文字负责}，否则训练会崩。这一招让平均分从 0.343 提到 0.431。\emph{凡是“模型序列里混入了非模型生成的外部内容”的 RL 场景，都该想到这一点。}
\item \textbf{简单的“结果奖励（RLVR）”就够用}：只用规则判答案对错（EM），不搞复杂过程奖励、不训神经奖励模型——既省事又稳、还不易被钻空子。
\item \textbf{多轮交错（推理↔检索）+ 自发涌现的行为}：模型在 RL 中\emph{自己}学会了多次检索、反复验证、纠错——这些没有被显式教过，是奖励信号逼出来的。
\item \textbf{算法选型有取舍}：GRPO 快但易崩、PPO 稳但慢，最终性能相当，求稳选 PPO；检索段数、组大小都存在“质量 vs. 稳定”的权衡。
\end{enumerate}

\textbf{对 AI 应用开发实习生的实际启示：}
\begin{itemize}[leftmargin=1.4em]
\item \emph{做检索类/问答类 Agent，别只停在“RAG 检一次”}。如果你的场景是多跳、需要分步查资料（如客服、研究助理、Deep-Research 类产品），Search-R1 这种“多轮自主检索”的形态是更强的目标形态。
\item \emph{如果你要用 RL 训这类 Agent，务必对“工具/检索返回的内容做掩码”}。这是从踩坑里换来的稳定性经验，能省掉大量“训练莫名崩掉”的调试时间。
\item \emph{奖励能用规则就别用模型}。在答案可自动判对错的任务上（数学、事实问答、代码能跑通），优先用可验证奖励，既便宜又稳。
\item \emph{先用结构化标签把“想/搜/答”框起来}（\tg{think}/\tg{search}/\tg{information}/\tg{answer}），系统层据此解析并接管搜索调用——这套“标签协议 + 系统接管工具执行”的工程范式，即使你暂时不做 RL、只做提示工程，也很值得直接借鉴。
\item \emph{基座可以从 Base 模型起步}。RL 能把 Base 和 Instruct 的差距补平，不一定非得用昂贵的指令模型当起点。
\end{itemize}

%==================================================================
\section*{参考文献（仅列最关键几条，英文原样）}
\small
\begin{enumerate}[leftmargin=2em]
\item Daya Guo, et al. DeepSeek-R1: Incentivizing reasoning capability in LLMs via reinforcement learning. arXiv preprint arXiv:2501.12948, 2025.
\item John Schulman, Filip Wolski, Prafulla Dhariwal, Alec Radford, Oleg Klimov. Proximal policy optimization algorithms. arXiv preprint arXiv:1707.06347, 2017.
\item Zhihong Shao, et al. DeepSeekMath: Pushing the limits of mathematical reasoning in open language models (GRPO). arXiv preprint arXiv:2402.03300, 2024.
\item Patrick Lewis, et al. Retrieval-augmented generation for knowledge-intensive NLP tasks (RAG). NeurIPS, 2020.
\item Harsh Trivedi, et al. Interleaving retrieval with chain-of-thought reasoning for knowledge-intensive multi-step questions (IRCoT). arXiv:2212.10509, 2022.
\item Shunyu Yao, et al. ReAct: Synergizing reasoning and acting in language models. ICLR, 2023.
\item Timo Schick, et al. Toolformer: Language models can teach themselves to use tools. NeurIPS, 2023.
\item Xiaoxi Li, et al. Search-o1: Agentic search-enhanced large reasoning models. arXiv:2501.05366, 2025.
\end{enumerate}

\end{document}
